\chapter{Diskussion}

Nach der Umsetzung und der Validierung werden einige Aspekte dieser Arbeit kritisch reflektiert.

\section{Reflektion des Projekts im Hinblick auf Zielsetzung}
\label{reflection}

Die Grundvoraussetzung, dass ein in EPK modellierter Prozess auf einem IoT-System ausgeführt und orchestiert werden kann ist mit 
dieser Arbeit belegt worden. Dieses Ziel wurde also erreicht. Ein kritischer Punkt, der nun reinkommt ist, 
dass nun Personal für diese Features geschult werden muss bzw. die Systeme, welche dieses Feature nutzen möchten dementsprechende Updates 
durchführen müssen. Auch kommt noch die Umrüstung der Industrieanlagen dazu, wo jede Anlage mit einem ESP32 verbunden werden muss, um eine 
Anbindung und Steuerung über bspw. ARIS zu gewährleisten.

\section{Gefahren für die Gültigkeit}

Auch wenn der Protoyp gemäß der Zielsetzung funktioniert, gibt es einige Risiken zu beachten. Besonders
im Sicherheitsaspekt.

\subsection{Absicherung des Animationssockets}

Zurzeit laufen die TCP-Sockets zur Animation während der Ausführung ohne jegliche Absicherung. Dies kann dazu 
führen, dass Daten abgelesen werden können zwischen dem ARIS-Server und Client. Hier kann durch Nutzung von TLS
die Verbindung abgesichert werden.

\subsection{Absicherung des Resource Directory Servers}

Auch beim Resource Direcotry Server ist es so, dass dieser auf Port 5683 hört, welcher nicht abgesichert ist. Dieser 
ist kritischer, da dort die Zieladdressen der ESP32 liegen, welche angesteuert werden sollen. Hier kann man eine Absicherung 
gewährleisten durch Nutzung von dTLS für den Resource Directory Server.

\section{Kritische Bewertung der Architektur und Systemlimitationen}

Die Arbeit wurde mit der Software ARIS umgesetzt, welche modifiziert wurde. Auch gab es hier einige kritische Punkte, während 
der Umsetzung welches diskutiert werden sollten.

\subsection{Latenz durch Netzwerkkommunikation}

Das Ziel der Arbeit war, dass Prozesse in Echtzeit von ARIS angesteuert werden. Durch Nutzung von CoAP auf Basis von UDP kann man 
einiges an Latenz und Overhead durch die kompakte Paketgröße einsparen. Jedoch wird es immer eine Latenz geben, welche abhängig von 
den Netzwerkgegebenheiten ist.

\subsection{Bottleneck durch die JVM auf Server und Client}

Neben der Netzwerklatenz kommt auch noch der Aspekt, dass die JVM durch ihre Garbage-Collection kein deterministisches Verhalten liefert und 
es zu unbestimmten Zeiten zu Verlangsamungen kommen kann. Dies kann auch das angestrebte Echtzeitverhalten der Orchestierung beeinflussen. In solchen
Fällen ist es hilfreich mit Sprachen wie C/C++ zu arbeiten, wo man durch eigene Speicherverwaltung eine stabile Performance gewährleisten kann.

\subsection{Fehlerbehandlung (synchron und asynchron)}

Aktuell ist die Fehlerbehandlung auf dem ESP32 durch ein passendes Request-Response Muster gewährleistet. Auf ARIS hingegen ist die Fehlerbehandlung mit
Laufzeitausnahmen umgesetzt. Hier kann durch normale Ausnahmen eine bessere Fehlerbehandlung umgesetzt werden, sodass es nicht zu einem hängendem Programm kommt.

\subsection{Continous Integration und Continous Delivery}

Die Nutzung von Continous Integration und Continous Delivery für ARIS ist durch deren Infrastruktur auf GitHub gegeben. Für die Firmware des ESP32 kann Continous-Integration 
umgesetzt werden durch Testen und Bauen des Codes auf Plattformen wie GitHub mit Actions oder GitLab mit GitLab CI. 

\section{Ausblick}

Neben der Umsetzung der Lösungen der kritischen Punkte, gibt es auch weitere mögliche Aussichten für das Resultat dieser Arbeit.

\subsection{Modularisierung durch Einführung von unterschiedlichen Notationsfrontends, Zielbackends und Eventloggingformaten}

Aktuell unterstützt der Protoyp die EPK-Notation, läuft auf ARIS und generiert Code für den ESP32. Jedoch können auch Frontends für andere Notationen 
wie BPMN oder BPMN-2.0 entwickelt werden und auch kann man für unterschiedliche Prozessverwaltungsplattformen das passende Backend entwickeln.
Neben diesen zwei Aspekten kann auch die Codegenerierung für weitere Mikrocontroller wie Arduino oder STM32 ausgeweitet werden. Auch kann das Logging der 
Ausführung in anderen Formaten umgesetzt werden, sodass weitere Auswertung auf Basis dieser Daten durchgeführt werden können.

\subsection{Open-Sourcing des Ausführungsmodells, Codegeneration, Event-Logging und Live-Monitoring}

Der Eigenanteil dieser Arbeit ist der Compiler, das Ausführungsmodell sowie die Codegenerierung und das Live-Monitoring. Dadurch, dass es sehr viel Potenzial gibt das ganze weiter 
auszubauen ist es sinnvoll, diese Anteile (mit Ausnahme der ARIS-Komponenten) der Öffentlichkeit zugänglich zu machen, um die Automatisierung und Orchestierung von industriellen Geschäftsprozessen 
durch Open-Source Entwicklung voranzutreiben.

\section{Fazit zur Praxistauglichkeit}

Durch Execution-Engines wie Flowable und Camunda hat man schon gesehen, dass Prozesse automatisiert werden können. Jedoch tut diese Arbeit auch neben dieser Prozessautomatisierung, ein Rahmenwerk festlegen, 
sodass Software gemäß von Geschäftsprozessen entwickelt werden kann. Diese prozessorientierte Softwareentwicklung kann in der Praxis sehr sinnvoll sein und durch Vereinfachung von Entwicklungsprozessen durch 
Lösungen wie Codegenerierung ist es umsetzbar. Kritisch kann es jedoch werden bei der Aufrüstung der industriellen Komponenten, da hier Anbindungen gemacht werden müssen zwischen der Anlage und dem Mikrocontroller, 
für eine Verbindung zwischen ARIS und den Anlagen wie in \ref{reflection} angesprochen. Es ist also im Fazit umsetzbar softwareseitig, jedoch kann es auf der Hardwareseite zu Schwierigkeiten kommen. Sei es durch Inkompatibilität 
oder einfach der fehlenden Möglichkeit der Anbindung an die Anlage. Es benötigt Anlagen mit Periphärien oder einer Schnittstelle, um einen Mikrocontroller zu integrieren. Hierbei kann man sich ein Beispiel an 
Periphärien wie CAN aus der Automobilindsutrie nehmen, welche CAN zur Realisierung von OBD verwenden.